% Options for packages loaded elsewhere
\PassOptionsToPackage{unicode}{hyperref}
\PassOptionsToPackage{hyphens}{url}
%
\documentclass[
]{article}
\usepackage{amsmath,amssymb}
\usepackage{iftex}
\ifPDFTeX
  \usepackage[T1]{fontenc}
  \usepackage[utf8]{inputenc}
  \usepackage{textcomp} % provide euro and other symbols
\else % if luatex or xetex
  \usepackage{unicode-math} % this also loads fontspec
  \defaultfontfeatures{Scale=MatchLowercase}
  \defaultfontfeatures[\rmfamily]{Ligatures=TeX,Scale=1}
\fi
\usepackage{lmodern}
\ifPDFTeX\else
  % xetex/luatex font selection
\fi
% Use upquote if available, for straight quotes in verbatim environments
\IfFileExists{upquote.sty}{\usepackage{upquote}}{}
\IfFileExists{microtype.sty}{% use microtype if available
  \usepackage[]{microtype}
  \UseMicrotypeSet[protrusion]{basicmath} % disable protrusion for tt fonts
}{}
\makeatletter
\@ifundefined{KOMAClassName}{% if non-KOMA class
  \IfFileExists{parskip.sty}{%
    \usepackage{parskip}
  }{% else
    \setlength{\parindent}{0pt}
    \setlength{\parskip}{6pt plus 2pt minus 1pt}}
}{% if KOMA class
  \KOMAoptions{parskip=half}}
\makeatother
\usepackage{xcolor}
\usepackage[margin=1in]{geometry}
\usepackage{graphicx}
\makeatletter
\newsavebox\pandoc@box
\newcommand*\pandocbounded[1]{% scales image to fit in text height/width
  \sbox\pandoc@box{#1}%
  \Gscale@div\@tempa{\textheight}{\dimexpr\ht\pandoc@box+\dp\pandoc@box\relax}%
  \Gscale@div\@tempb{\linewidth}{\wd\pandoc@box}%
  \ifdim\@tempb\p@<\@tempa\p@\let\@tempa\@tempb\fi% select the smaller of both
  \ifdim\@tempa\p@<\p@\scalebox{\@tempa}{\usebox\pandoc@box}%
  \else\usebox{\pandoc@box}%
  \fi%
}
% Set default figure placement to htbp
\def\fps@figure{htbp}
\makeatother
\setlength{\emergencystretch}{3em} % prevent overfull lines
\providecommand{\tightlist}{%
  \setlength{\itemsep}{0pt}\setlength{\parskip}{0pt}}
\setcounter{secnumdepth}{-\maxdimen} % remove section numbering
\usepackage[spanish]{babel}
\usepackage{amsmath}
\usepackage{fontspec}
\usepackage{unicode-math}
\usepackage{graphicx}
\usepackage{adjustbox}
\usepackage{tikz}
\usepackage{pgfplots}
\usetikzlibrary{3d,babel}
\usepackage{bookmark}
\IfFileExists{xurl.sty}{\usepackage{xurl}}{} % add URL line breaks if available
\urlstyle{same}
\hypersetup{
  hidelinks,
  pdfcreator={LaTeX via pandoc}}

\author{}
\date{\vspace{-2.5em}}

\begin{document}

\section{Question}\label{question}

En una comunidad se realizó un estudio sobre el número de hijos por
familia. Los resultados obtenidos se presentan en la siguiente tabla:

\includegraphics[width=12cm,height=\textheight,keepaspectratio]{tabla_frecuencias.png}

Con base en la información de la tabla, ¿cuál de las siguientes
afirmaciones sobre la frecuencia relativa de las familias con 3 hijos es
correcta?

\subsection{Answerlist}\label{answerlist}

\begin{itemize}
\tightlist
\item
  El 72.1\% de las familias tienen 3 hijos
\item
  El 27.9\% de las familias tienen 3 hijos
\item
  El 24.3\% de las familias tienen 3 hijos
\item
  El 73.0\% de las familias tienen 3 hijos
\end{itemize}

\section{Solution}\label{solution}

Para interpretar correctamente la frecuencia relativa, debemos recordar
que:

\begin{enumerate}
\def\labelenumi{\arabic{enumi}.}
\tightlist
\item
  \textbf{Frecuencia relativa} = Frecuencia absoluta / Total de datos
\item
  La frecuencia relativa representa la proporción de datos en cada
  categoría
\item
  Para expresarla como porcentaje, multiplicamos por 100
\end{enumerate}

\textbf{Análisis de la tabla:}

\begin{itemize}
\tightlist
\item
  Total de familias encuestadas: 111
\item
  Familias con 3 hijos: 31
\item
  Frecuencia relativa: 0.279
\end{itemize}

\textbf{Cálculo del porcentaje:}

\[\text{Porcentaje} = \text{Frecuencia relativa} \times 100 = 0.279 \times 100 = 28\%\]

Por lo tanto, la interpretación correcta es: \textbf{El 27.9\% de las
familias tienen 3 hijos}

\textbf{Análisis de los distractores comunes:}

\begin{itemize}
\tightlist
\item
  Confundir frecuencia absoluta con porcentaje (usar directamente el
  número de familias como porcentaje)
\item
  No convertir la frecuencia relativa a porcentaje (dejar el valor
  decimal)
\item
  Calcular el porcentaje sobre una categoría diferente
\item
  Calcular el porcentaje complementario (100\% - porcentaje correcto)
\end{itemize}

Es fundamental distinguir entre: - \textbf{Frecuencia absoluta}:
cantidad de casos - \textbf{Frecuencia relativa}: proporción respecto al
total (valor entre 0 y 1) - \textbf{Porcentaje}: frecuencia relativa ×
100

\subsection{Answerlist}\label{answerlist-1}

\begin{itemize}
\tightlist
\item
  Falso. Esta no es la interpretación correcta de los datos.
\item
  Verdadero. Esta es la interpretación correcta de la frecuencia
  relativa convertida a porcentaje.
\item
  Falso. Esta no es la interpretación correcta de los datos.
\item
  Falso. Esta no es la interpretación correcta de los datos.
\end{itemize}

\section{Meta-information}\label{meta-information}

exname: Interpretación de frecuencia relativa en tabla de distribución
extype: schoice exsolution: 0100 exshuffle: TRUE exsection:
Estadística/Tablas de frecuencia/Interpretación
exextra{[}icfes\_competencia{]}: interpretacion\_representacion
exextra{[}icfes\_componente{]}: aleatorio exextra{[}icfes\_nivel{]}: 2
exextra{[}icfes\_contexto{]}: comunitario

\end{document}
